\chapter{绪论}
\setcounter{chapter}{1}
\chapterintro{
\textbf{【教学大纲要求】}

\imp{掌握：}职业卫生与职业医学的概念和任务；劳动与健康的关系；工作条件（生产工艺过程、工作过程、工作环境）；职业性有害因素及其来源；职业病（广义与立法意义的概念）、工作有关疾病、工伤的涵义；职业性病损的致病模式\keyword{*}（高危人群的概念）；三级预防原则与目的；职业病的特点；职业病的诊断原则。

\imp{熟悉：}现阶段职业卫生的特点和发展趋势（职业有害因素范围扩展和职业危害转嫁；产业结构与就业状态变化与农民工职业卫生问题；职业紧张和心理障碍；职业伤害与职业卫生突发事件；职业安全、职业卫生和环境保护的融合；新理论和技术在职业卫生中的应用；接触纳米材料的职业卫生问题；普及基本职业卫生服务，加强一级预防和人类工效学研究）。

（注：\keyword{*} 标示教学难点）
}

% ----- 1.1 基本概念 -----
\section{职业卫生与职业医学概述}

\begin{defbox}
\textbf{职业卫生（Occupational Health）}：以职业人群为主要对象，研究劳动条件对健康的影响，主要任务是识别、评价、预测、控制不良劳动条件，为制定措施改善劳动条件、保护劳动者健康、提高作业能力提供科学依据。关注群体，着眼预防（主要属于一级预防）。

\textbf{职业医学（Occupational Medicine）}：以劳动者为对象，对受到职业危害因素损害或存在潜在健康危险的个体，采用临床医学的技术与方法，对发生的职业病、职业相关疾病和早期健康损害进行检测、诊断、治疗和康复处理。
\end{defbox}

\subsection{职业、工作与劳动}

\begin{defbox}
\textbf{职业}：人们在社会中所从事的作为主要生活来源的具有特定专业和技能要求的工作类别。特点：专业性、群体性（稳定性和发展性）。

\textbf{工作}：个人在某个工作环境中所承担的具体任务或活动。特点：具体性、任务导向性（多样性和灵活性）。

\textbf{劳动}：人类为了创造物质、精神财富而进行的有目的的活动，是人类生存和发展的基础。特点：基础性、体力和脑力结合、创造性。
\end{defbox}

\begin{keybox}
\textbf{劳动与人的关系——双重性：}
\begin{enumerate}[leftmargin=*]
    \item \textbf{社会性与功利性}：劳动者生存所必需；体现人身价值，对幸福感和社会身份产生重要影响
    \item \textbf{健康相关性}：维持健康所必需；职业有关因素导致疾病；造成安全事故；污染环境影响健康
\end{enumerate}
\end{keybox}

% ----- 1.2 职业有害因素 -----
\section{职业有害因素及其来源}

\begin{defbox}
\textbf{职业有害因素（Occupational Hazards）}：又称职业病危害因素、职业性有害因素。在职业场所中产生或存在的、可能对职业人群健康、安全和作业能力造成不良影响的因素或条件，包括化学、物理、生物等因素。
\end{defbox}

\begin{keybox}
\textbf{职业有害因素的分类（按来源）：}

\textbf{一、生产过程中的有害因素}
\begin{enumerate}[leftmargin=*]
    \item \imp{化学因素}：
    \begin{itemize}
        \item 生产性毒物：铅、汞、CO$_2$、香蕉水、农药等
        \item 生产性粉尘：矽尘、石棉尘、煤尘、有机粉尘等
    \end{itemize}
    \item \imp{物理因素}：
    \begin{itemize}
        \item 异常气象条件：高温、低温、高湿、高气压、低气压
        \item 噪声、振动
        \item 辐射：可见光、紫外线、红外线、微波、激光、X射线、γ射线
    \end{itemize}
    \item \imp{生物因素}：炭疽、布氏病、森林脑炎、过敏性肺炎、艾滋
\end{enumerate}

\textbf{二、劳动过程中的有害因素}
\begin{itemize}[leftmargin=*]
    \item 劳动组织和制度不合理（劳动时间过长、休息制度不合理）
    \item 劳动强度过大
    \item 精神（心理）紧张
    \item 器官或系统过度紧张（视力紧张、发音器官紧张）
    \item 不良体位（下背痛、颈肩腕损伤、下肢静脉曲张）
    \item 不合理工具设备
\end{itemize}

\textbf{三、生产环境中的有害因素}
\begin{itemize}[leftmargin=*]
    \item 自然环境因素（夏季太阳辐射、冬季寒冷、高山高原环境）
    \item 工作场所设计不合理（厂房低矮狭窄、布局不合理）
    \item 卫生技术设施不足（通风换气、照明、防尘防毒防噪防振设备）
\end{itemize}

\textbf{国际上按因素性质分类（5类）：}
\begin{enumerate}[leftmargin=*]
    \item 化学因素：粉尘、毒性气体、金属、有机溶剂等
    \item 物理因素：高温、振动、噪声、异常气压、辐射等
    \item 生物因素：细菌、病毒、真菌、立克次体等
    \item 人机工效学因素：生物力学、人体尺寸、工具、人机界面等
    \item 社会心理因素：紧张、压力等
\end{enumerate}
\end{keybox}

% ----- 1.3 职业病与工作有关疾病 -----
\section{职业病与工作有关疾病}

\subsection{职业病}

\begin{defbox}
\textbf{职业病（Occupational Disease）}：劳动者在职业活动中，因接触粉尘、放射性物质和其他有毒、有害因素而引起的疾病。具有\keyword{医学和法律双重含义}。

\textbf{法定职业病的4个基本条件：}
\begin{enumerate}[leftmargin=*]
    \item 患病主体：必须是企业、事业单位或个体经济组织的劳动者
    \item 患病环境：必须在从事职业活动的过程中产生
    \item 致病原因：必须因接触粉尘、放射性物质和其他有毒、有害物质等职业病危害因素引起
    \item 疾病类型：必须是国家公布的职业病分类和目录所列的职业病
\end{enumerate}
\end{defbox}

\begin{keybox}
\textbf{职业病的特点：}
\begin{enumerate}[leftmargin=*]
    \item \imp{病因明确}——可以预防
    \item \imp{明确的剂量-反应关系}
    \item \imp{接触人群中常有一定的发病率}（集群性）
    \item 如\imp{早期诊断、及时治疗}，预后良好
    \item \imp{多数无特效治疗方法}
\end{enumerate}

\textbf{我国职业病分类目录：}2024版，12大类135种。

\textbf{制定职业病清单主要考虑的因素：}职业病的严重程度；职业病的社会与经济影响；国家经济承受能力；国际上可供参考的惯例。
\end{keybox}

\subsection{工作有关疾病}

\begin{defbox}
\textbf{工作有关疾病（Work-Related Disease）}：与多因素相关的疾病，在职业活动中，由于职业性有害因素等多种因素的作用，导致劳动者罹患某种疾病或潜在疾病显露或原有疾病加重。
\begin{itemize}[leftmargin=*]
    \item 职业因素是疾病发生和发展的诸多因素之一（非唯一病因）
    \item 职业因素影响健康，从而促使潜在的病症显露或加重
    \item 通过改善工作条件，可使所患疾病得到控制或缓解
\end{itemize}

\textbf{常见的工作有关疾病：}
\begin{itemize}[leftmargin=*]
    \item 行为与身心疾病：焦虑、抑郁、神经衰弱
    \item 非特异性呼吸道疾病：慢性支气管炎、哮喘
    \item 其他：高血压、消化性溃疡、腰背痛
\end{itemize}
\end{defbox}

\begin{defbox}
\textbf{职业特征（Occupational Stigma）}：职业性有害因素导致的生理性代偿，如胼胝、皮肤色素增加等。不属于疾病范畴。

\textbf{工伤（Injuries）}：在生产劳动过程中，由于操作、设备、管理等原因和不可预测的偶然因素所造成的工人身体伤害、残疾甚至死亡。
\end{defbox}

% ----- 1.4 职业病诊断与三级预防 -----
\section{职业病的诊断与处理}

\begin{keybox}
\textbf{职业病诊断原则：}
\begin{enumerate}[leftmargin=*]
    \item 详细的职业接触史
    \item 相应的临床表现和辅助检查
    \item 职业卫生现场调查资料
    \item 综合分析，排除其他病因所致类似疾病
\end{enumerate}

\textbf{处理原则：}
\begin{enumerate}[leftmargin=*]
    \item 积极有效的治疗
    \item 调离作业环境
    \item 职业伤害鉴定及落实待遇
    \item 报告制度：急性中毒24小时报告，慢性中毒15天内报告；急性中毒同时发生3名死亡，电话上报；职业炭疽1人以上时，电话上报
\end{enumerate}
\end{keybox}

\subsection{三级预防原则}

\begin{keybox}
\textbf{职业病三级预防原则：}

\textbf{一级预防（病因预防）：}
\begin{itemize}[leftmargin=*]
    \item 以无毒代有毒，改进生产工艺
    \item 制定容许接触水平
    \item 发现职业禁忌证
    \item 开展健康教育
\end{itemize}

\textbf{二级预防（发病预防）：}
\begin{itemize}[leftmargin=*]
    \item 职业健康监护
    \item 早期检测损害
    \item 早期诊断
    \item 早期治疗处理
\end{itemize}

\textbf{三级预防（对症康复预防）：}
\begin{itemize}[leftmargin=*]
    \item 脱离接触
    \item 积极治疗
    \item 促进康复
    \item 防止并发症
\end{itemize}
\end{keybox}

% ----- 1.5 职业卫生服务 -----
\section{职业卫生服务与职业健康监护}

\begin{defbox}
\textbf{职业健康监护}：以预防为目的，根据劳动者职业接触史，通过定期或不定期的医学健康检查和健康相关资料收集，连续地监测劳动者健康状况，分析健康变化与所接触的职业病危害因素的关系，并及时将结果报告给用人单位和劳动者本人，以便采取干预措施。

\textbf{职业卫生服务（Occupational Health Service, OHS）}：由具有预防职能的服务机构，负责向用人单位和劳动者提出建议，帮助用人单位为劳动者创造一个安全与健康的工作环境，以促进劳动者的体力与脑力健康，使工作适合于劳动者的生理特点。

\textbf{基本职业卫生服务（BOHS）}：以预防和控制职业病、保护和促进劳动者的健康和工作能力为目的，采用科学合理和社会可接受的方法，通过初级卫生保健方式，为所有劳动者提供的必要的职业卫生技术服务。
\end{defbox}

\begin{tipbox}
\textbf{综合性职业卫生服务的发展方向：}
\begin{itemize}[leftmargin=*]
    \item 预防疾病→增进健康：由预防严重职业病转向更为积极地增进劳动者健康
    \item 单一监控→综合服务：由单一监测、监护、评价转为监测-监护-评价-控制配套系统服务
    \item 重点人群→所有劳动者：由对部分职业人群监测转为对所有劳动者的监测、监护
\end{itemize}
\end{tipbox}

\section{职业卫生防治工作的主要内容}

\begin{enumerate}[leftmargin=*]
    \item \imp{职业卫生监督}：预防性卫生监督（新、改、扩建项目）；经常性卫生监督（现有项目）
    \item \imp{职业流行病学的应用}：
    \begin{itemize}
        \item 阐明职业性损害在人群中的分布、发生和发展规律
        \item 发现职业性有害因素对健康的影响
        \item 为制定、修订职业卫生标准和职业病诊断标准提供依据
        \item 评价职业卫生和职业病防治工作质量及其预防措施的效果
    \end{itemize}
    \item \imp{职业流行病学中的重要概念——健康工人效应（HWE）}：
    \begin{defbox}
    \textbf{健康工人效应（Healthy Worker Effect, HWE）}：在进行职业流行病学调查时，观察到暴露于职业有害因素的工人总死亡率低于一般人群总死亡率的现象。产生原因：
    \begin{enumerate}[leftmargin=*]
        \item \textbf{初次就业选择效应}：健康状况较好者更可能进入就业队伍
        \item \textbf{继续就业选择效应}：患病或健康状况不佳者可能被解雇或主动离职
        \item \textbf{时间相关效应}：健康工人效应在就业初期最明显，随工龄延长而减弱
    \end{enumerate}
    HWE的存在可能导致低估职业有害因素的危害程度，是职业流行病学研究中的重要偏倚来源。
    \end{defbox}
    \item \imp{培训与宣传}：新的职业危害因素；危害防护的知识与技能；有关法律、法规、标准、制度
\end{enumerate}

\section{现阶段性职业卫生挑战}

\begin{enumerate}[leftmargin=*]
    \item 职业有害因素范围扩展与职业危害转移
    \item 产业结构与就业状态变化与农民工职业健康问题
    \item 职业紧张和心身疾病
    \item 职业伤害与职业卫生突发事件
    \item 职业安全、职业卫生和环境保护的内在联系
    \item 新理论和技术在职业卫生中的应用
    \item 从接触标志物到职业健康管理
    \item 普及基本职业卫生服务，加强一级预防，加强工效学研究
\end{enumerate}

\section{新业态劳动者的职业健康问题}

\begin{keybox}
随着互联网经济的快速发展，网约配送员（外卖骑手）、网约车司机、网课教师等新业态劳动者大量涌现。这些职业具有工作时间灵活、工作场所不固定、劳动关系复杂等特点，面临着与传统职业不同的职业健康风险。

\textbf{一、网约配送员（外卖骑手）}
\begin{itemize}[leftmargin=*]
    \item \imp{主要职业相关健康风险：}
    \begin{enumerate}[leftmargin=*]
        \item \textbf{职业伤害}：交通事故（最主要风险，时间压力导致超速、违章）；跌倒、碰撞
        \item \textbf{肌肉骨骼疾病}：长期骑行→腰背痛、颈椎病；搬运重物→腰椎劳损
        \item \textbf{不良气象条件}：夏季高温→中暑风险；冬季寒冷→冻伤、关节疼痛
        \item \textbf{心理紧张}：时效考核压力、差评焦虑→职业紧张、睡眠障碍
        \item \textbf{接触空气污染物}：道路PM$_{2.5}$、汽车尾气等
        \item \textbf{不良生活方式}：饮食不规律→消化系统疾病；憋尿→泌尿系统疾病
    \end{enumerate}
    \item \imp{预防建议：}
    \begin{enumerate}[leftmargin=*]
        \item 优化配送算法和时效考核（避免不合理的时效压力）
        \item 交通安全教育和防护装备（头盔、反光背心）
        \item 高温季节提供防暑降温保障
        \item 职业伤害保险全覆盖
        \item 定期健康检查
        \item 合理安排工作时间，保障基本休息权
    \end{enumerate}
\end{itemize}

\textbf{二、网约车司机}
\begin{itemize}[leftmargin=*]
    \item \imp{主要职业相关健康风险：}
    \begin{enumerate}[leftmargin=*]
        \item \textbf{肌肉骨骼疾病}：长期久坐→腰背痛、颈肩腕损伤、下肢静脉曲张
        \item \textbf{振动危害}：全身振动→腰椎间盘突出、脊柱退行性改变
        \item \textbf{视觉紧张}：长时间注视道路和屏幕→视力疲劳、干眼
        \item \textbf{心理紧张}：收入压力、交通拥堵、乘客纠纷→职业紧张、焦虑
        \item \textbf{排尿系统问题}：如厕不便→憋尿→泌尿系统感染、膀胱功能障碍
        \item \textbf{饮食和作息不规律}：胃肠道疾病、肥胖、代谢综合征
        \item \textbf{接触空气污染}：汽车尾气暴露→呼吸系统疾病风险
    \end{enumerate}
    \item \imp{预防建议：}
    \begin{enumerate}[leftmargin=*]
        \item 合理规划驾驶时间（连续驾驶不超过2小时）
        \item 车辆减振设计和座椅人体工学改进
        \item 定期停车休息、伸展活动
        \item 健康教育和职业健康监护
        \item 心理支持和社会支持体系建设
    \end{enumerate}
\end{itemize}

\textbf{三、网课教师}
\begin{itemize}[leftmargin=*]
    \item \imp{主要职业相关健康风险：}
    \begin{enumerate}[leftmargin=*]
        \item \textbf{发音器官过度紧张}：长时间连续授课→声带小结、慢性咽炎、声音嘶哑
        \item \textbf{视觉紧张（VDT作业）}：长时间注视屏幕→视力下降、干眼症、视疲劳
        \item \textbf{静力作业和不良体位}：长期久坐→颈肩腕综合征、腰背痛
        \item \textbf{职业紧张（信息性劳动—综合式）}：课程质量要求高、学生互动压力→焦虑、职业倦怠（Burnout）
        \item \textbf{辐射暴露}：长期使用电脑→微波和极低频电磁场暴露
        \item \textbf{作息不规律}：直播课程时间不固定→昼夜节律紊乱、睡眠障碍
    \end{enumerate}
    \item \imp{预防建议：}
    \begin{enumerate}[leftmargin=*]
        \item 控制连续授课时间（每45--60分钟休息一次）
        \item 科学用嗓训练和护嗓（多饮水、避免清嗓子）
        \item 符合工效学的工作站设计（可调节桌椅、屏幕高度与眼齐平）
        \item 合理安排课程时间，保障规律作息
        \item 职业紧张干预（心理支持、团队建设）
        \item 定期健康检查和嗓音功能评估
    \end{enumerate}
\end{itemize}

\begin{exambox}
\textbf{新业态劳动者职业健康问题的共同特点：}
\begin{enumerate}[leftmargin=*]
    \item 劳动关系复杂化（平台化用工），传统职业卫生监管体系覆盖不足
    \item 工作场所分散化、工作时间碎片化，增加了职业健康管理的难度
    \item 社会保障和职业病保障体系不健全
    \item 预防措施需要政府、平台、从业者三方共同努力
\end{enumerate}
\end{exambox}
\end{keybox}

% ----- 1.6 复习题 -----
\section{复习题}

\begin{enumerate}[leftmargin=*, itemsep=8pt]
    \item \textbf{【名词解释】}职业卫生；职业医学；职业有害因素
    \item \textbf{【名词解释】}职业病；工作有关疾病；职业特征；工伤
    \item \textbf{【名词解释】}职业健康监护；基本职业卫生服务（BOHS）；健康工人效应（HWE）
    \item \textbf{【简答题】}简述职业病的特点（5条）。
    \item \textbf{【简答题】}简述法定职业病的4个基本条件。
    \item \textbf{【简答题】}简述职业病的三级预防原则及各级的主要内容。
    \item \textbf{【简答题】}简述职业有害因素的分类（按来源分为生产过程、劳动过程和生产环境三类）。
    \item \textbf{【简答题】}简述职业病的诊断原则和处理原则。
    \item \textbf{【简答题】}简述综合性职业卫生服务的三个发展方向。
    \item \textbf{【思考题】}网约配送员（外卖骑手）、网约车司机、网课教师可能面临哪些职业健康风险？如何预防？
\end{enumerate}

\subsection*{参考答案要点}

\begin{enumerate}[leftmargin=*, itemsep=5pt]
    \item \textbf{职业卫生}：以职业人群为主要对象，研究劳动条件对健康的影响，主要任务是识别、评价、预测、控制不良劳动条件，保护劳动者健康。\textbf{职业医学}：以劳动者为对象，对受到职业危害因素损害的个体进行检测、诊断、治疗和康复处理。\textbf{职业有害因素}：在职业场所中产生或存在的、可能对职业人群健康、安全和作业能力造成不良影响的因素或条件。
    \item \textbf{职业病}：劳动者在职业活动中因接触粉尘、放射性物质和其他有毒、有害因素而引起的疾病，具有医学和法律双重含义。\textbf{工作有关疾病}：职业因素为诸多致病因素之一的疾病。\textbf{职业特征}：职业性有害因素导致的生理性代偿（如胼胝）。\textbf{工伤}：生产劳动过程中因操作、设备、管理等原因造成的身体伤害、残疾甚至死亡。
    \item \textbf{职业健康监护}：以预防为目的，通过医学健康检查和健康相关资料收集，连续监测劳动者健康状况。\textbf{BOHS}：以预防和控制职业病为目的，通过初级卫生保健方式为所有劳动者提供的必要职业卫生技术服务。\textbf{健康工人效应}：暴露于职业有害因素的工人总死亡率低于一般人群的现象，由初次就业选择和继续就业选择效应产生，是职业流行病学的重要偏倚来源。
    \item 五大特点：病因明确可预防；明确的剂量-反应关系；接触人群中常有一定发病率（集群性）；早期诊断及时治疗预后良好；多数无特效治疗方法。
    \item 四个基本条件：患病主体为企业/事业单位/个体经济组织的劳动者；患病环境为从事职业活动过程中；致病原因为接触职业病危害因素；疾病类型必须为国家公布的职业病目录所列疾病。
    \item 一级预防（病因预防）：无毒代毒、制定容许接触水平、发现禁忌证、健康教育。二级预防（发病预防）：健康监护、早期发现、早期诊断、早期治疗。三级预防（对症康复）：脱离接触、积极治疗、促进康复、防止并发症。
    \item 生产过程：化学因素（毒物、粉尘）、物理因素（异常气象、噪声振动、辐射）、生物因素（炭疽、布氏杆菌等）。劳动过程：劳动组织不合理、强度过大、精神紧张、器官过度紧张、不良体位、不合理工具。生产环境：自然环境因素、场所设计不合理、卫生技术设施不足。
    \item 诊断原则：详细职业接触史+相应临床表现+职业卫生现场调查+综合分析排除其他病因。处理原则：积极治疗+调离作业环境+职业伤害鉴定落实待遇+及时报告。
    \item 三个发展方向：预防疾病→增进健康；单一监控→综合服务（监测-监护-评价-控制配套）；重点人群→所有劳动者。
\end{enumerate}

\newpage

% =====================================================================
%  CHAPTER 2: 职业生理学、心理学与工效学
% =====================================================================
\newpage